\subsection{Compile deadlines and repeated build costs, 15 November}
\label{sec:compile-cost}
Texloom runs installed TeX engines as child processes. Rust prepares a
private copy of the project, checks dependencies, runs the required passes
and seals the resulting PDF. The old build service stopped a build after
120 seconds, limited a TeX pass to 90 seconds and bounded bibliography
generation. The new path lets these processes finish or be cancelled.
Version and lookup probes keep their own deadlines. The compiler still has
six passes to settle its references and generated inputs, and the process
memory limit remains 2 GiB.

A synthetic Lua\LaTeX{} fixture completed a first pass of
\textbf{121.83 seconds} and an eligible build in \textbf{122.97 seconds}.
The auxiliary flag prevents another delay on the second pass. The process
tests cover explicit cancellation and descendant cleanup in a run with
no deadline. This control exercises the real engine through the native
build service.

The release diagnostic used MiKTeX pdf\LaTeX{}, Xe\LaTeX{} and Lua\LaTeX{}
on the same Windows host. For each engine, the workload used a plain
article, an auxiliary cross-reference and a project with 24 MiB of unrelated
source. Every build created a fresh source snapshot. The first baseline
plain-article medians were 1.10, 1.82 and 1.94 seconds, respectively.
The two engine passes accounted for roughly 0.76, 1.43 and 1.57 seconds.
Version identification alone took 0.26--0.29 seconds. These intervals exclude
editor debounce, WebView transport and visible presentation.

The revised path checks all dependencies together after recording
them, and resolves each repeated recorder path once per pass. A repeat
build may reuse engine identification for matching compiler bytes, alias,
project and environment. Its successful passes must still report the
expected runtime and format. Snapshot isolation, TeX execution,
dependency checks and convergence remain part of each build.

\begin{table}[htbp]
\small\centering
\begin{tabular}{@{}lrrr@{}}
\toprule
Plain article & Fresh check & Reused check & Paired saving\\
\midrule
pdfLaTeX & 1074 & 782 & 280\\
XeLaTeX & 1911 & 1702 & 294\\
LuaLaTeX & 2005 & 1715 & 325\\
\bottomrule
\end{tabular}
\caption{Alternating fresh and reused version checks, in milliseconds.
Each engine has five measured pairs after warm-up. The last column is the
median of within-pair differences.}
\label{tab:compile-pairstats}
\end{table}

All 90 comparisons against the baseline and 45 fresh/reused pairs retained
identical extracted text, full single-page rasters, pass counts, arguments
and used-project-input lists. These repeated synthetic conditions share
one machine. Background workload and filesystem cache state varied.
The protocol and raw-result identities are bound by
\path{docs/evidence/compile-performance-2026-11-15.json}.

The first pass accounts for much of the wait. The service then runs a second
pass to verify stability in the measured cases. A provisional preview at the
first successful pass could show content sooner while the verified result
continues to settle. Reusing auxiliary state is a separate experiment
that needs cold-build equivalence controls for references, changed files,
bibliographies and dynamic TeX. Established continuous builders such as
latexmk manage the same class of dependencies and reruns~\cite{latexmkManual}.

A separate Node transport diagnostic measured 267 ms median to parse,
validate and convert a JSON reply carrying 8 MiB of PDF-pattern bytes.
The JSON payload was 29.6 MB. It measures frontend CPU work before PDF
rendering. Tauri's binary response API supports a follow-up that can
avoid byte-array JSON and retain revision-bound build
metadata~\cite{tauriBinaryIPC}. A packaged Windows measurement would assess
delivery and presentation together.

\subsection{Binary delivery development follow-up}

The desktop command now encodes a versioned envelope containing a UTF-8
JSON metadata header and the original PDF bytes. The browser checks the
envelope and then applies the existing build identity, revision, mapping
and eligibility rules. A typed byte view avoids parsing and validating
each PDF byte as a JSON number. The PDF surface still copies the view
before handing it to PDF.js, whose worker may transfer buffer ownership.
Native service callers retain their original result type. Tauri's
postMessage fallback can return the envelope as a number array; that
slower path validates each byte before decoding the same envelope.

\begin{table}[htbp]
\small\centering
\begin{tabular}{@{}lrrr@{}}
\toprule
Pattern buffer & Legacy decode & Binary decode & Paired saving\\
\midrule
1 MiB & 126.04 & 1.00 & 124.97\\
8 MiB & 895.15 & 6.07 & 889.38\\
\bottomrule
\end{tabular}
\caption{Node component diagnostic, milliseconds. Each condition has five
alternating measured pairs after one warm-up pair. These are decoding
costs, not end-to-end compile or visible-preview times.}
\label{tab:binary-compile-reply}
\end{table}

The legacy arm loads the preserved validator from commit
\code{0c3cd89}; the binary arm uses the new decoder and validator. Both
intervals include the surface's byte copy, with exact-byte checks after
timing. The 8 MiB envelope occupies 8,388,812 bytes, compared with
29,644,012 JSON bytes. The earlier 267 ms observation and two interim
development blocks remain separate; uncontrolled machine state and JIT
execution can affect these small repeated samples. No inference about
whole-application speed follows from their ratio.

Generated Rust envelope bytes cross the actual ordered-client validator
in regression checks. A real PDF fixture retains identical extracted text,
geometry and rendered pixels after delivery. Other controls cover failed
builds, malformed envelopes, Unicode metadata, cancellation, candidate
revisions, PDF size limits and fallback decoding. The receipt is
\path{docs/evidence/binary-compile-reply-2026-11-15.json}. Native WebView
dispatch and physical presentation time remain unmeasured.

\subsection{First-pass observation feasibility, 16 November}

A direct-engine control retained first-pass PDFs before later passes
overwrote them. Across pdf\LaTeX{}, Xe\LaTeX{} and Lua\LaTeX{}, plain
articles retained the same text and pixels, while references and contents
changed before settling. Each deliberately failed second pass followed a
readable first PDF. These 12 synthetic conditions establish why an early
PDF needs a provisional lifecycle; they bypass the application build service.

An opt-in observer now exposes first-pass bytes from the supervised build
library after recorder discovery and frozen-input/dependency checks. Its
type is explicitly unverified, carries diagnostics, and provides no receipt
or mapping authority. Reads are bounded at 60 MiB and require a PDF header;
consumers must still parse before presentation. Cancellation is checked
around the callback. Ordinary desktop builds continue to return the settled
result and perform no early byte read.

The production-service corpus exercises seven conditions under each engine:
plain article, references, contents, later-pass failure, first-pass compile
error, callback cancellation and changed frozen input after observation.
All 21 controls pass. The nine successful final PDFs retain the baseline's
text, character geometry and full-page 96 dpi rasters. Three plain first
PDFs match the final output; the six reference/contents first PDFs differ.
The later-failure controls remain failed and ineligible. Cancellation stops
before a second completed pass, and changed frozen input prevents a final
build while the original project remains unchanged.

The final native run primes successful compiler identification with the
no-observer baseline, then records reuse explicitly for observed builds.
Its first-available and final-return times are descriptive endpoints, with
one observation per engine/condition. Those durations have different probe
conditions and do not estimate callback overhead or application speedup.
Earlier direct and native development runs remain separate. The protocol is
\path{research/FIRST-PASS-PREVIEW-PLAN.md}; the source and result identities
are in \path{docs/evidence/first-pass-preview-2026-11-16.json}.

The subsequent desktop experiment used a separate read-only view keyed to the
project, revision and request. It retained the last verified PDF after failure.
The owner reported that the tested candidate at \code{0bf8f8c} worked well and
asked to defer early presentation. Normal Texloom builds now wait for complete
compilation before showing a new PDF. The source retains the experimental path
in a disabled configuration. The owner feedback is qualitative, and the corpus
provides no input-to-display measurements.

\subsection{Auxiliary reuse and clean-build equivalence}

A constructed case tests whether an unchanged auxiliary state guarantees the
same output as a clean build. The old source defines a macro as \texttt{B} and
writes that definition to its auxiliary file. The revised source supplies
\texttt{A} only when the definition is absent, then writes the value it used.
The two revised builds use identical source. One begins with an empty output
directory. The other starts with the old auxiliary file.

For each of pdf\LaTeX{}, Xe\LaTeX{} and Lua\LaTeX{}, the clean revision prints
\texttt{Branch A} and the seeded revision prints \texttt{Branch B}. Two passes
within each arm produce matching auxiliary hashes, extracted text and full-page
96 DPI rasters. Across arms, the text and rasters differ. The old auxiliary
file is another input, and a document can intentionally use such persistent
state. A reuse policy that promises clean-build equivalence must account for
that distinction even after convergence. The experiment consists of one
constructed case per engine and provides no estimate of how often projects
exhibit this behavior or how much time caching would save. The protocol, logs and
per-pass outputs are retained with the development evidence.

\subsection{Precompiled preamble feasibility}

The \texttt{mylatexformat} driver can dump a document's preamble into a format
file so that the compiler can start with its packages already loaded~\cite{mylatexformat}.
The experiment uses MiKTeX pdf\LaTeX{} on a plain article and a synthetic
article that loads TikZ libraries, amsmath, geometry and hyperref. Each format
comes from a seed document with the same preamble and a different body. Both
arms build the same revised source twice in fresh directories with empty
auxiliary state. One warm-up pair precedes five measured pairs per fixture,
with alternating arm order. Timings cover the two compiler processes, and
format generation is measured separately.

\begin{table}[htbp]
\small\centering
\begin{tabular}{@{}lrrrr@{}}
\toprule
Fixture & Standard & Preloaded & Paired saving & Dump\\
\midrule
Plain article & 743 & 717 & 44 & 318\\
TikZ article & 1722 & 852 & 830 & 907\\
\bottomrule
\end{tabular}
\caption{Preamble feasibility, milliseconds. Build times and paired savings
are medians over five measured pairs per fixture. Dump gives the one-time
format generation interval.}
\label{tab:preamble-feasibility}
\end{table}

All ten measured pairs match extracted text, character geometry, full-page
96 DPI rasters and link rectangles and destinations. The initial SyncTeX
command used a relative source path, and the installed client found no tag
in either arm. A separate check of the retained files uses their recorded
absolute source paths. It matches the queried forward geometry and reverse
source line in all ten measured pairs, without rerunning the timing trials.

The format files occupy 13.4 and 16.9 MB. In the TikZ fixture, the final-pass
recorder lists 131 distinct input paths in the standard arm and 20 in the
preloaded arm. Package inputs consumed while dumping the format are therefore
part of the dependency inventory even if a later pass omits them. In the
invalidation control, a changed preamble macro prints \texttt{Revised} in a
standard build and \texttt{Original} with the stale format. Any production
reuse must bind preamble and dependency identity across both stages. These
small repeated observations describe one installed engine and two fixtures.
They exclude the desktop pipeline and do not establish broader package,
engine or SyncTeX equivalence. Texloom's normal builds use the standard
compiler format.

\subsection{Engine identification across project switches}

The earlier optimization remembered one compiler identification. When a user
switched projects or compiler commands, a successful build replaced that
entry. An in-memory cache now retains up to eight recently successful project
and command pairs. A match requires the same executable path, fingerprint and
version-affecting environment. The build consumes the entry and returns it to
the cache only after success. It executes every TeX pass and checks the
runtime banner, dependencies and build outputs.

The paired native workload alternates two independent synthetic projects
under each MiKTeX engine. It compares retained identification with a fresh
version probe, changing arm order by trial. Each project has one warm-up pair
and five measured pairs per engine. The medians of the paired savings for the
two projects are 151 and 326 ms for pdf\LaTeX{}, 361 and 286 ms for
Xe\LaTeX{}, and 254 and 324 ms for Lua\LaTeX{}. All 36 pairs, including the
warm-ups, match text, character geometry, 96 DPI rasters, pass count, engine
arguments and recorded dependency identities. The source bytes remain
unchanged. The 30 measured pairs represent repeated native runs on one host.
These intervals exclude editor work, IPC and PDF presentation. The aggregate
observations and evidence identities are recorded in
\path{docs/evidence/compile-followup-2026-11-18.json}.
